\section{论文工作安排计划}
\label{section-plan}

\subsection{工作进度安排}
\label{section-plan-arrangement}

表\ref{tab:schedule}为本工作的工作进度安排表。

\begin{table}[htbp]
  \centering
  \begin{tabular}{|c|l|}
    \hline
    \textbf{时间} & \multicolumn{1}{|c|}{\textbf{任务}} \\
    \hline
    2025.09 -- 2025.10 & 状态机需求与模型数据集整理 \\
    \hline
    2025.11 -- 2026.02 & 多步式建模方法实现与验证 \\
    \hline
    2026.03 -- 2026.06 & 验证场景与性质生成方法开发 \\
    \hline
    2026.07 -- 2026.10 & 基于剖面的状态机验证方法构建 \\
    \hline
    2026.11 -- 2027.01 & 模型迭代修复方法实现及闭环测试 \\
    \hline
    2027.02 -- 2027.04 & 博士论文撰写与答辩准备 \\
    \hline
  \end{tabular}
  \caption{任务时间表}
  \label{tab:schedule}
\end{table}


\subsection{论文工作基础}
\label{section-plan-foundation}

本研究立足于控制系统功能安全设计与验证的实际需求，已在多个关键方面积累了扎实的工作基础，为后续研究的顺利开展提供了有力支撑。

在数据集准备方面，本研究已收集并整理了来自多个控制系统的功能安全需求数据集。这些数据来源于公开数据集、工具自带案例以及工业实践中的典型需求，涵盖了不同规模和复杂度的控制系统。具体数据如表\ref{tab:combined_stats}所示，共包含9个系统的101条功能安全需求，这些需求以结构化条件判断、状态转移和流程控制为核心，包含对控制系统实时性、交互性和可靠性相关的需求描述。例如，"如果当前状态为高风险，则将在250毫秒内检测到紧急情况"表达了控制系统的实时性需求；"如果启动自动控制按钮被按下，而安全带不可用，则发出报警"体现了控制系统的交互性需求；"当到达最终锚点时，速度降至0，PBA关闭"则体现了系统的可靠性需求。基于这些需求，本研究已构建了相应的状态机参考模型，模型中包含的状态数量从5到27不等，事件数量从3到13不等，变量数量从1到12不等，迁移数量从7到27不等，形成了涵盖不同规约复杂度的模型集合，为后续实验评估提供了丰富的数据基础。



\begin{table}[htbp]
\centering
\adjustbox{max width=\textwidth}{% 自动缩放至页面宽度
\begin{tabular}{|l|c|c|c|c|c|c|c|c|c|c|}
\hline
 & BSN\upcite{VOGEL2023107100} & CARA\upcite{7092662} & Elevator\upcite{Decan2020} & Microwave\upcite{Decan2020} & PBA\upcite{Feiler2012} & Radar & Stopwatch\upcite{Decan2020} & TCS\upcite{GDMS_TCS} & VHL\upcite{houdek2013system} & 合计 \\
\hline
需求 & 11 & 14 & 4 & 9 & 8 & 13 & 5 & 15 & 22 & 101 \\
\hline
状态 & 10 & 5 & 8 & 10 & 9 & 27 & 6 & 12 & 14 & 101 \\
\hline
事件 & 3 & 7 & 3 & 13 & 8 & 13 & 5 & 7 & 9 & 68 \\
\hline
变量 & 9 & 6 & 3 & 7 & 5 & 12 & 1 & 8 & 9 & 60 \\
\hline
迁移 & 7 & 12 & 10 & 17 & 9 & 25 & 7 & 27 & 27 & 141 \\
\hline
\end{tabular}}
\caption{各系统功能安全需求与状态机模型规约数据统计表}
\label{tab:combined_stats}
\end{table}

在技术工具与平台建设方面，本研究已有重要进展。我们设计并实现了一种基于领域特定语言（DSL）的状态机描述语言（pyfcstm），该语言已开源并提供使用\footnote{项目地址：https://github.com/HansBug/pyfcstm}。相较于现有的各种状态机元模型，我们的DSL设计更加注重对LLM的友好性，提供了清晰的语法结构和语义定义，同时为后续的形式化验证做了充分准备，支持各种形式化表达式系统。这一基础设施为研究内容一中多步式建模方法的实现提供了重要支撑。

在核心技术积累方面，我们对形式化验证工具UPPAAL的核心算法库进行了深入分析，结合其相关的算法讲解论文，已经掌握了大部分时间自动机的基础算法和操作，并将部分改进代码贡献合并到了官方开源库中。这一技术积累为研究内容三中基于验证剖面的状态机验证方法提供了坚实的算法基础，后续我们也将复用一部分来自UPPAAL的算法或者技术。

此外，针对实际工业应用中常见的非典型状态机验证需求（如伪状态、多层复合状态、带有历史状态的状态机等），我们正在开发一个专门的状态机算法库，旨在提供对这些复杂状态机结构的标准算法支持。该库的核心功能之一是将非典型状态机归约为简单状态机或时间自动机，以便进行形式化验证。这一工作将极大增强本研究对实际工业场景的适应能力。

本研究还得到了航天三院和17所相关项目的实际需求和数据支持，这些工业界的实际应用场景为研究方法提供了真实的验证环境和需求输入，确保研究成果具备实际应用价值。



\subsection{可能遇到的问题以及解决途径}
\label{section-plan-problems}

在本课题的研究过程中，预期可能会遇到以下几个关键问题：

第一，LLM生成的状态机模型在复杂控制系统场景下的准确性和完备性问题。由于控制系统往往涉及复杂的状态迁移逻辑、严格的时间约束和并发行为，LLM可能难以一次性生成完全符合要求的状态机模型，可能出现状态缺失、迁移条件不完整或时间属性表述错误等情况。针对这一问题，我们计划采取多步迭代的生成策略，将复杂的建模任务分解为多个子任务，通过引入模型自检和初步验证环节，对LLM生成的模型进行质量评估，并基于反馈进行迭代优化。同时，我们将建立控制系统领域知识库，为LLM提供领域特定的约束和模式，引导其生成更加准确和完备的模型。

第二，时间属性在状态机验证中的表达与处理复杂性。控制系统对实时性要求极高，时间属性的正确表达和验证是确保系统功能安全的关键。然而，时间属性的形式化表述和验证本身就具有很高的复杂性，加上LLM对时间概念的理解可能存在偏差，这会给验证过程带来挑战。为解决这一问题，我们将基于时间自动机理论，设计专门的时间状态机验证框架，明确定义时间属性的表达形式和验证规则。同时，我们将开发时间约束的规范化表示方法，减少LLM在处理时间概念时的歧义性，并通过示例学习和提示工程增强LLM对时间属性的理解能力。

第三，验证剖面生成的质量与覆盖度保证问题。验证剖面的质量直接影响到后续验证的效果，如果生成的验证剖面不能充分覆盖系统的关键行为和边界条件，就无法有效发现模型中的潜在缺陷。针对这一问题，我们将采用多源信息融合的策略，结合需求文档、模型结构信息和领域知识，生成多样化的验证剖面。同时，我们将引入覆盖度评估机制，对生成的验证剖面进行质量评估，并通过迭代优化不断提高其覆盖度。此外，我们还将研究基于模型敏感路径分析的剖面生成方法，优先覆盖那些最容易出现问题的关键路径。

第四，模型修复过程中可能引入新缺陷的问题。当LLM基于验证反馈对模型进行修复时，可能会在修复原有缺陷的同时引入新的问题，特别是对于那些存在复杂交互关系的模型元素。为了应对这一挑战，我们将建立修复影响分析机制，在应用修复方案前评估其可能产生的影响范围，避免引发连锁反应。同时，我们将采用保守的修复策略，优先选择最小化的修复方案，并通过回归验证确保修复后的模型不仅解决了已知问题，而且没有引入新的缺陷。

第五，状态空间爆炸问题在复杂控制系统验证中的挑战。尽管本研究采用了基于验证剖面的聚焦式验证策略，但对于特别复杂的控制系统，仍然可能面临状态空间爆炸的问题。为了缓解这一问题，我们将研究基于抽象精化的验证方法，通过合理的抽象减少状态空间规模，同时保持验证结果的可信度。此外，我们还将探索符号执行和偏序归约等优化技术，进一步提高验证效率。

通过上述解决方案，我们有信心能够克服研究过程中可能遇到的各种挑战，确保研究工作的顺利进行和预期目标的达成。
